\documentclass[11pt]{article}
\usepackage[margin=1in]{geometry}
\usepackage[T1]{fontenc}
\usepackage[utf8]{inputenc}
\usepackage{lmodern}
\usepackage{setspace}
\usepackage{amsmath, amssymb, amsthm}
\usepackage{booktabs}
\usepackage{multirow}
\usepackage{enumitem}
\usepackage{xcolor}
\usepackage{hyperref}
\usepackage{titlesec}
\usepackage{array}

\definecolor{titleblue}{RGB}{23,54,93}
\hypersetup{
    colorlinks=true,
    linkcolor=titleblue,
    urlcolor=titleblue,
    citecolor=titleblue
}
\titleformat{\section}{\large\bfseries\color{titleblue}}{\thesection.}{0.5em}{}
\titleformat{\subsection}{\normalsize\bfseries\color{titleblue}}{\thesubsection}{0.5em}{}
\setlist[itemize]{leftmargin=1.5em}
\setlist[enumerate]{leftmargin=1.6em}
\setstretch{1.12}

\title{\textbf{Methodology and Technical Design for}\\[0.2em]\textbf{Learning-Augmented Robust Traffic Engineering}\\[0.4em]
\large Detailed Design, Algorithms, Optimization, and Implementation Plan}
\author{ECE GY 7363: Communications Networks II -- Design and Algorithms}
\date{April 2026}

\begin{document}
\maketitle

\begin{abstract}
This document presents the detailed methodology for the proposed project on learning-augmented robust traffic engineering in software-defined networks. It explains the design formulation, optimization model, machine learning techniques, baseline algorithms, resilience evaluation strategy, implementation architecture, and experimental methodology. The project combines time-series traffic prediction with multi-commodity flow optimization to proactively route traffic in a capacitated network. The primary design goal is to minimize congestion while preserving feasibility and resilience under demand variations and link failures. The document is intended as the technical blueprint that connects the project idea to the actual implementation and evaluation pipeline.
\end{abstract}

\section{Project Goal}
The project studies how to improve traffic engineering decisions in a communication network when traffic demand varies over time and the network may experience failures. The central idea is to combine prediction and optimization in an end-to-end control loop:
\[
\text{historical traffic} \rightarrow \text{traffic prediction} \rightarrow \text{routing optimization} \rightarrow \text{resilience evaluation}.
\]

The key hypothesis is that a predictor trained on past traffic matrices can estimate near-future demand well enough to help the controller make better routing decisions than methods that react only to the current state. The project also studies when such prediction-enhanced routing remains effective under failures and when more conservative methods are preferable.

\section{System Model}
\subsection{Network Representation}
The communication network is represented as a directed graph
\[
G = (V,E),
\]
where:
\begin{itemize}
    \item $V$ is the set of nodes, representing routers or switches,
    \item $E$ is the set of directed links,
    \item each link $(i,j) \in E$ has capacity $c_{ij} > 0$.
\end{itemize}

The network operates over discrete time epochs $t = 1,2,\dots,T$. At each epoch, there is a traffic matrix that specifies the traffic demand between source-destination pairs.

\subsection{Traffic Demand Model}
Each source-destination pair is treated as a commodity $k \in \mathcal{K}$. Let:
\begin{itemize}
    \item $s_k$ denote the source of commodity $k$,
    \item $t_k$ denote the destination of commodity $k$,
    \item $d_k(t)$ denote the demand of commodity $k$ at time $t$.
\end{itemize}

The set of all commodity demands at time $t$ forms the traffic matrix. The ML component will predict a one-step-ahead estimate $\hat{d}_k(t+1)$ using historical traffic observations.

\section{Design Formulation}
\subsection{Decision Variables}
The routing design is defined by continuous flow variables:
\[
f_{ij}^{k}(t) \ge 0,
\]
where $f_{ij}^{k}(t)$ is the amount of traffic of commodity $k$ sent on directed link $(i,j)$ at time $t$.

This means the design problem is not simply selecting one path per flow. Instead, it determines how much traffic of each commodity should traverse each link so that all constraints are respected and the routing objective is optimized. This is the standard multi-commodity flow formulation used in traffic engineering.

\subsection{Flow Conservation Constraints}
For every commodity $k$ and node $v \in V$, flow must satisfy:
\[
\sum_{(u,v)\in E} f_{uv}^{k}(t) - \sum_{(v,w)\in E} f_{vw}^{k}(t) =
\begin{cases}
-d_k(t), & v = s_k, \\
d_k(t), & v = t_k, \\
0, & \text{otherwise.}
\end{cases}
\]

These constraints ensure:
\begin{itemize}
    \item the source injects the required traffic,
    \item the destination absorbs the required traffic,
    \item every intermediate node preserves flow balance.
\end{itemize}

\subsection{Capacity Constraints}
Each link has limited capacity, so the total routed traffic on a link must not exceed that link's capacity:
\[
\sum_{k \in \mathcal{K}} f_{ij}^{k}(t) \le c_{ij}, \qquad \forall (i,j)\in E.
\]

This is a core design constraint because it captures congestion and makes routing a capacity-sharing problem across all commodities.

\subsection{Nonnegativity Constraints}
All routed flows must be nonnegative:
\[
f_{ij}^{k}(t) \ge 0, \qquad \forall (i,j)\in E, \; k\in\mathcal{K}.
\]

\subsection{Objective Function}
The primary design objective is to minimize the worst congestion level in the network. Define an auxiliary variable $U(t)$ representing the maximum link utilization at time $t$. Then impose:
\[
\sum_{k \in \mathcal{K}} f_{ij}^{k}(t) \le U(t)\, c_{ij}, \qquad \forall (i,j)\in E.
\]

The optimization objective is:
\[
\min U(t).
\]

This is a classical min-max congestion objective. It attempts to spread traffic in a balanced way so that no single link becomes a bottleneck.

\subsection{Why This Formulation Fits the Project}
This design formulation is well suited to the project because:
\begin{itemize}
    \item it directly models concurrent traffic demands,
    \item it respects physical link capacities,
    \item it produces routing decisions that can be compared across different demand and failure scenarios,
    \item it integrates naturally with predicted demand from the ML model.
\end{itemize}

\section{Optimization Techniques}
\subsection{Primary Optimization Technique: Linear Programming}
The main optimization method used in the project is \textbf{linear programming (LP)}. The problem is linear because:
\begin{itemize}
    \item the decision variables are continuous flow values,
    \item flow conservation constraints are linear,
    \item capacity constraints are linear,
    \item the min-max utilization objective becomes linear after introducing $U(t)$.
\end{itemize}

The LP will be solved at each time epoch using the current observed demand or the predicted next-step demand, depending on the algorithm being tested.

\subsection{Learning-Augmented Optimization}
The project uses \textbf{learning-augmented optimization}. In this framework:
\begin{itemize}
    \item an ML model predicts the next traffic matrix,
    \item the predicted demand is treated as input to the optimization model,
    \item the optimizer computes routes based on predicted rather than only current demand.
\end{itemize}

This is not a black-box ML controller. The optimization layer still guarantees capacity-feasible routing. ML only improves the quality of the input information used by the solver.

\subsection{Rolling-Horizon Control}
The project will operate in a \textbf{rolling-horizon} manner. At each time step:
\begin{enumerate}
    \item collect the most recent traffic history,
    \item predict the next traffic matrix,
    \item solve the routing LP using the predicted matrix,
    \item apply or evaluate the routing decision,
    \item observe the realized traffic and repeat.
\end{enumerate}

This reflects how an SDN controller would operate in practice.

\subsection{Optional Extensions}
If time permits, the following advanced optimization extensions may be considered:
\begin{itemize}
    \item \textbf{Robust optimization:} hedge against demand uncertainty or predefined failure scenarios.
    \item \textbf{Fairness-aware optimization:} add constraints or secondary objectives related to equitable flow treatment.
    \item \textbf{Mixed-integer formulations:} enforce path-based or discrete-routing decisions. This is lower priority because it is significantly harder computationally.
\end{itemize}

\section{Algorithms and Techniques}
\subsection{Baseline Routing Algorithms}
The project will compare the proposed approach against several baseline methods.

\paragraph{Shortest-Path Routing.}
Each demand is routed on a shortest path using Dijkstra's algorithm. This baseline is simple, interpretable, and widely used as a reference point. However, it ignores shared congestion effects between commodities.

\paragraph{Optimization with Current Demand Only.}
At each time step, solve the LP using the current observed traffic matrix rather than a prediction. This baseline isolates the value of optimization without ML.

\paragraph{Conservative Non-Predictive Routing.}
Inflate or smooth the observed demand and solve the LP on that conservative estimate. This serves as a basic robust baseline.

\subsection{Traffic Prediction Algorithms}
The project will use a staged comparison of forecasting methods.

\paragraph{Moving Average.}
Predict the next traffic value as the average of the previous few time steps. This is easy to implement and provides a low-complexity benchmark.

\paragraph{Linear Regression.}
Use past traffic entries or aggregated temporal features to fit a simple linear predictor. This gives a stronger statistical baseline than moving average.

\paragraph{Multi-Layer Perceptron (Optional).}
Flatten a short history of traffic matrices into a feature vector and predict the next matrix using a feed-forward neural network. This captures nonlinear patterns but does not explicitly model sequence memory.

\paragraph{LSTM.}
The main ML method will be a Long Short-Term Memory network. LSTM is appropriate because traffic matrices are sequential data with short-term temporal correlation, daily trends, and occasional bursts. LSTM can learn these dependencies more effectively than static models.

\paragraph{GRU (Optional).}
If needed, a Gated Recurrent Unit model may be used as a lighter recurrent alternative.

\subsection{Priority of Advanced ML Models}
To strengthen the project, three more modern ML directions can be considered in addition to LSTM. Their priority for this project is as follows.

\paragraph{Priority 1: Transformer-based time-series model.}
This is the highest-priority modern extension. A Transformer-based predictor is highly relevant because traffic matrices form a multivariate time series, and Transformer architectures are effective at modeling longer temporal dependencies than standard recurrent models. In the context of communication networks, this is useful when traffic exhibits periodicity, delayed correlations, or burst patterns that span many time steps. It is also a practical extension because it can be added without changing the optimization layer of the project.

\paragraph{Priority 2: Graph Neural Network (GNN) or Spatio-Temporal GNN.}
This is the second-priority modern extension and is highly relevant to networks because the underlying system is naturally represented as a graph. A GNN can incorporate topology information directly, allowing the predictor to model how traffic correlations depend on the physical or logical network structure. This is useful when nearby or strongly connected nodes influence each other's traffic patterns. For network-focused research value, GNN-based models are more structurally meaningful than plain sequence models, but they are somewhat more complex to implement and tune.

\paragraph{Priority 3: Graph Transformer or Spatio-Temporal Graph Transformer.}
This is the most ambitious and research-oriented option. A graph transformer combines temporal attention with graph structure learning, making it theoretically the most expressive of the three advanced choices. In communication networks, it is especially relevant when both long-range temporal dependencies and topological interactions matter simultaneously. However, it has the highest implementation complexity, the largest tuning burden, and the greatest risk for a semester project. For that reason, it is best treated as a stretch goal rather than the default model.

\paragraph{Recommended model hierarchy.}
For a strong but manageable project, the recommended hierarchy is:
\begin{itemize}
    \item \textbf{Base model:} LSTM,
    \item \textbf{Best modern extension:} Transformer,
    \item \textbf{Most network-structured extension:} GNN,
    \item \textbf{Stretch model:} Graph Transformer.
\end{itemize}

\paragraph{Interpretation for this project.}
If the goal is the safest path to a high-quality result, the project should implement LSTM first and add a Transformer as the main ``latest'' comparison model. If the goal is to emphasize that the problem is inherently graph-structured, then a GNN-based predictor is the most network-relevant advanced model. If the goal is maximum novelty and ambition, then a graph transformer is the most advanced option, but it should only be attempted after the rest of the pipeline is stable.

\subsection{Failure Analysis Techniques}
The project will study resilience using:
\begin{itemize}
    \item \textbf{Random single-link failures,} to test average resilience,
    \item \textbf{Critical-link failures,} selected based on utilization, betweenness centrality, or topological importance,
    \item \textbf{Traffic spike plus failure scenarios,} to test combined stress conditions.
\end{itemize}

\subsection{Performance Evaluation Techniques}
Performance will be compared using:
\begin{itemize}
    \item maximum link utilization,
    \item average link utilization,
    \item number of overloaded links,
    \item fraction of demand satisfied,
    \item average path length,
    \item Jain's fairness index,
    \item degradation under failure,
    \item ML prediction error using MAE, RMSE, and relative error.
\end{itemize}

\section{Detailed End-to-End Algorithm}
The overall algorithmic pipeline is as follows.

\begin{enumerate}
    \item Load the network topology and capacities.
    \item Generate or load a sequence of traffic matrices.
    \item Split traffic data into training, validation, and test periods.
    \item Train the traffic prediction model on historical data.
    \item For each test time step:
    \begin{enumerate}[label=\alph*.]
        \item extract the recent traffic history window,
        \item predict the next traffic matrix,
        \item solve the routing LP using the predicted demand,
        \item evaluate the produced routing under realized demand,
        \item simulate link-failure scenarios and record performance metrics.
    \end{enumerate}
    \item Repeat the same evaluation loop for all baseline methods.
    \item Compare methods statistically and graphically.
\end{enumerate}

\section{Implementation Architecture}
\subsection{Planned Software Stack}
The project will be implemented in Python. The core stack is:
\begin{itemize}
    \item \texttt{networkx} for graph and topology handling,
    \item \texttt{numpy} and \texttt{pandas} for data processing,
    \item \texttt{PyTorch} for ML model development,
    \item \texttt{cvxpy} or \texttt{pulp} for optimization,
    \item \texttt{matplotlib} and \texttt{seaborn} for plots and comparisons.
\end{itemize}

\subsection{Planned Code Modules}
The codebase will be organized into the following modules:
\begin{enumerate}
    \item \textbf{topology.py}: load and store graph structure, capacities, and metadata.
    \item \textbf{traffic\_generator.py}: create synthetic time-varying traffic matrices.
    \item \textbf{predictors.py}: moving average, linear regression, and LSTM models.
    \item \textbf{optimizer.py}: LP-based routing solver.
    \item \textbf{baselines.py}: shortest path and non-predictive routing methods.
    \item \textbf{failure\_simulator.py}: link-removal scenarios and resilience evaluation.
    \item \textbf{metrics.py}: congestion, fairness, and forecasting metrics.
    \item \textbf{run\_experiments.py}: end-to-end execution pipeline.
    \item \textbf{plots.py}: figures and comparison visualizations.
\end{enumerate}

\subsection{Data Flow in the Implementation}
The runtime flow will be:
\begin{enumerate}
    \item topology is loaded,
    \item traffic matrices are generated or loaded,
    \item training data is used to fit the predictor,
    \item prediction output is passed to the optimizer,
    \item optimizer output is passed to the evaluator,
    \item results are stored and plotted for all methods.
\end{enumerate}

\section{Experimental Design}
\subsection{Network Instances}
The experiments will use one or more standard backbone-style networks such as NSFNET, Abilene, or GEANT. Synthetic random topologies may also be included to test generality.

\subsection{Traffic Generation Strategy}
When real traces are unavailable, traffic matrices will be generated using:
\begin{itemize}
    \item base demand matrix,
    \item periodic modulation,
    \item Gaussian or bounded random noise,
    \item occasional burst events,
    \item hotspot pairs with elevated traffic.
\end{itemize}

This creates realistic and reproducible demand variation while maintaining control over difficulty.

\subsection{Train-Test Setup}
The traffic sequence will be divided into:
\begin{itemize}
    \item training set for learning predictor parameters,
    \item validation set for model selection,
    \item test set for final routing comparisons.
\end{itemize}

\subsection{Failure Scenario Design}
Failure testing will include:
\begin{itemize}
    \item one random failed link per episode,
    \item highest-utilization link failure,
    \item top betweenness-centrality link failure,
    \item multiple independent random failure samples.
\end{itemize}

\section{Complexity and Practical Considerations}
The LP routing model is polynomial-time solvable but may grow in size with the number of commodities and links. Therefore:
\begin{itemize}
    \item initial experiments will use medium-sized topologies,
    \item commodity count may be restricted if necessary,
    \item path-based simplifications can be considered if runtime becomes large.
\end{itemize}

The LSTM adds training cost, but prediction inference during evaluation is typically cheap. This makes the approach suitable for a repeated control loop once the predictor is trained.

\section{Expected Findings}
The expected outcomes are:
\begin{itemize}
    \item optimization-based routing should outperform shortest-path routing under congestion,
    \item prediction-enhanced routing should outperform current-demand-only routing when traffic exhibits temporal structure,
    \item under severe prediction error or extreme failures, conservative baselines may remain competitive,
    \item resilience analysis will show whether predictive routing improves or harms robustness under failure.
\end{itemize}

\section{Main Risks and Mitigation}
\begin{itemize}
    \item \textbf{Risk:} optimization becomes slow on large topologies. \\
    \textbf{Mitigation:} begin with medium-size graphs and scale only after validation.
    \item \textbf{Risk:} synthetic traffic is too noisy for ML to help. \\
    \textbf{Mitigation:} generate traffic with controlled temporal structure and evaluate forecasting accuracy separately.
    \item \textbf{Risk:} full robust optimization is too ambitious. \\
    \textbf{Mitigation:} prioritize scenario-based resilience evaluation first.
\end{itemize}

\section{Why the Technical Plan is Strong}
This methodology is strong academically because it combines:
\begin{itemize}
    \item a mathematically rigorous design formulation,
    \item a classical and interpretable optimization framework,
    \item a modern ML enhancement,
    \item systematic baseline comparisons,
    \item explicit resilience and robustness evaluation,
    \item a realistic implementation plan.
\end{itemize}

That combination makes the project well aligned with the course and well positioned for a high-quality final report and presentation.

\section{Conclusion}
This document provides the full technical plan for implementing the proposed project. The design formulation is a dynamic multi-commodity flow model on a capacitated graph. The main optimization technique is linear programming with a min-max utilization objective. The main ML technique is LSTM-based traffic prediction, compared against classical forecasting baselines. The project will evaluate routing quality, resilience to failures, and the value of learning-augmented optimization through a structured experimental pipeline. As a result, this methodology serves as a complete blueprint for carrying the project from idea to implementation and evaluation.

\end{document}
